\subsubsection{Prufer code}

É uma maneira de codificar uma árvore rotulada em uma sequência de números
de maneira única. É uma bijeção entre as árvores geradoras de um grafo completo
e as sequências numéricas de $n-2$ inteiros no intervalo $[0, n-1]$.

Ajuda a resolver problemas combinatóricos.

Aqui, é mostrado o número de maneiras de adicionar arestas a um grafo de modo
que ele fique conexo. Não são consideradas árvores de um único vértice.

\textbf{Construção:} selecionamos a folha da árvore com o menor número. Remove
ele da árvore. Escreve o número do vértice que tava conectado à folha. Depois
de $n-2$ iterações, sobrarão $2$ vértices e o algoritmo acaba. Aqui um jeito
de fazer isso linearmente.

\inccode{grafos/prufer_code_construction.cpp}

\textbf{Propriedades:}

\begin{enumerate}
    \item Depois de construir, dois vértices faltam. Um deles é o maior $n-1$. Mas
    nada pode ser dito sobre o outro.
    \item Todo vértice aparece no código de Prufer o grau dele menos uma vezes.
\end{enumerate}

Dá pra restaurar a árvore a partir do código de Prufer linearmente

\inccode{grafos/prufer_code_restoration.cpp}

\textbf{Fórmula de Cayley:} O número de árvores geradoras num grafo completo é
trivialmente $n^{n-2}$ pelo código de Prufer.

\textbf{Número de maneiras de tornar um grafo conexo:} Dado um grafo de $k$ componentes,
quero o número de maneiras de adicionar $k-1$ arestas de modo que o grafo fique conexo.
É como se tivéssemos um grafo de $k$ componentes e quiséssemos aplicar a fórmula de Cayley.
Sejam $d_i$ os graus. Temos que $\sum\limits_{i=1}^k d_i = 2k-2$. Isso quer dizer que o vértice
$i$ aparece $d_i-1$ vezes no código de Prufer. O código de Prufer de uma árvore de $k$ vértices
tem tamanho $k-2$, então o número de maneiras de escolher um código com $k-2$ números em que
o número $i$ aparece exatamente $d_i-1$ vezes é igual ao coeficiente multinomial
$$ \binom{k-2}{d_1 - 1, \dots, d_k-1}$$
O fato de que cada aresta adjacente ao vértice $i$ multiplica a resposta por $s_i$ (a quantidade de
caras na componente) implica que a resposta é
$$ \sum\limits_{d_i \geq 1, \, \sum d_i = 2k-2} s_1^{d_1} \dots s_k^{d_k} \binom{k-2}{d_1 - 1, \dots, d_k-1}$$
Trivialmente, isso é
$$ s_1 \dots s_k (s_1 + \dots + s_k)^{k-2} = s_1 \dots s_k n^{k-2}$$
Acidentalmente essa fórmula também funciona para $k=1$.